%%%% APÊNDICE A

\chapter{Contrato de Comunicação ARM--FPGA}\label{cap:apendicea}

Este apêndice consolida o contrato de comunicação utilizado pelo protótipo. Ele complementa a descrição do \autoref{cap:proposta} e deve permanecer alinhado aos arquivos \path{cpp/src/fpga_shared_stream.h}, \path{cpp/src/fast_receiver.cpp}, \path{vhdl/arm_fpga_shared_stream_bridge.vhd} e \path{docs/arm-fpga-shared-memory-stream.md}.

\section{Registradores}\label{sec:apendiceRegistradores}

A \autoref{tab:apendiceRegistradores} consolida os registradores de identificação, controle das filas e telemetria. Os campos de desempenho são produzidos pelo motor e lidos pelo ARM; escrever 1 em \codigo{PERF_CTRL} reinicia seus contadores.

\begin{table}[htb]
\caption{Registradores da interface MMIO}
\label{tab:apendiceRegistradores}
\begin{tabularx}{\textwidth}{llllX}
\toprule
\textbf{Offset} & \textbf{Nome} & \textbf{Acesso} & \textbf{Dono} & \textbf{Descrição} \\ \midrule
\codigo{0x000} & \codigo{MAGIC} & RO & FPGA & Valor \codigo{0x48465431}. \\
\codigo{0x004} & \codigo{VERSION} & RO & FPGA & Versão do protocolo. \\
\codigo{0x008} & \codigo{CTRL} & RW & ARM & Bit 0 reinicia ponteiros. \\
\codigo{0x00C} & \codigo{STATUS} & RO & FPGA & Bits de disponibilidade das filas. \\
\codigo{0x010} & \codigo{TX_HEAD} & RW & ARM & Publicação de comandos. \\
\codigo{0x014} & \codigo{TX_TAIL} & RO & FPGA & Consumo de comandos. \\
\codigo{0x018} & \codigo{RX_HEAD} & RO & FPGA & Publicação de respostas. \\
\codigo{0x01C} & \codigo{RX_TAIL} & RW & ARM & Consumo de respostas. \\
\codigo{0x020} & \codigo{TX_DEPTH} & RO & FPGA & Profundidade da fila TX. \\
\codigo{0x024} & \codigo{RX_DEPTH} & RO & FPGA & Profundidade da fila RX. \\
\codigo{0x028} & \codigo{SLOT_WORDS} & RO & FPGA & Palavras por quadro. \\
\codigo{0x030} & \codigo{PERF_CTRL} & RW & ARM & Bit 0 reinicia contadores. \\
\codigo{0x034} & \codigo{PERF_CLOCK_HZ} & RO & FPGA & Frequência usada na conversão de ciclos. \\
\codigo{0x038} & \codigo{PERF_COUNT} & RO & FPGA & Quantidade de eventos medidos. \\
\codigo{0x03C} & \codigo{PERF_LAST_LATENCY} & RO & FPGA & Latência mais recente, em ciclos. \\
\codigo{0x040} & \codigo{PERF_MIN_LATENCY} & RO & FPGA & Menor latência, em ciclos. \\
\codigo{0x044} & \codigo{PERF_MAX_LATENCY} & RO & FPGA & Maior latência, em ciclos. \\
\codigo{0x048--0x04C} & \codigo{PERF_SUM_LATENCY} & RO & FPGA & Soma de latências em 64 bits. \\
\codigo{0x050} & \codigo{PERF_CMD_STALL} & RO & FPGA & Ciclos de espera no comando. \\
\codigo{0x054} & \codigo{PERF_RSP_STALL} & RO & FPGA & Ciclos de espera na resposta. \\ \bottomrule
\end{tabularx}
\fonte{Elaborado pelos autores conforme \codigo{fpga_shared_stream.h} e \codigo{arm_fpga_shared_stream_bridge.vhd}}
\end{table}

\section{Regras de filas}\label{sec:apendiceFilas}

As filas TX e RX usam a mesma semântica:

\begin{itemize}
  \item vazia: \codigo{head == tail};
  \item cheia: \codigo{next(head) == tail};
  \item capacidade útil: \codigo{DEPTH - 1};
  \item produtor escreve payload antes de publicar \codigo{head};
  \item consumidor lê payload antes de publicar \codigo{tail}.
\end{itemize}

Para \codigo{DEPTH = 64} e \codigo{SLOT_WORDS = 8}, a base TX é \codigo{0x100}, cada slot tem 32 bytes e a base RX é \codigo{0x900}.

\section{Formato de evento}\label{sec:apendiceEvento}

O evento enviado do ARM para o FPGA usa oito palavras. A palavra de sequência é preservada para correlação entre entrada e saída. O identificador de símbolo é numérico para evitar processamento textual em hardware. O preço é escalado por $10^4$ para representar casas decimais com inteiros.

A \autoref{tab:apendiceEvento} apresenta a ordem das palavras. A ordem é parte do protocolo: alterar uma posição exige atualizar software, RTL, testes e documentação.

\begin{table}[htb]
\caption{Formato do quadro ARM--FPGA}
\label{tab:apendiceEvento}
\begin{tabularx}{\textwidth}{llX}
\toprule
\textbf{Palavra} & \textbf{Campo} & \textbf{Semântica} \\ \midrule
\codigo{word0} & Sequência & Número de sequência da mensagem sintética. \\
\codigo{word1} & Símbolo & Identificador numérico atribuído pelo receptor. \\
\codigo{word2} & Preço & Preço em escala $10^4$. \\
\codigo{word3} & Quantidade & Quantidade agregada no nível de preço. \\
\codigo{word4} & Tipo de evento & Atualização/inserção, remoção ou reset. \\
\codigo{word5} & Lado & Compra ou venda. \\
\codigo{word6} & Reserva & Campo reservado para extensão futura. \\
\codigo{word7} & Reserva & Campo reservado para extensão futura. \\ \bottomrule
\end{tabularx}
\fonte{Elaborado pelos autores conforme o contrato implementado em C++ e VHDL}
\end{table}

\section{Formato de resposta}\label{sec:apendiceResposta}

O software interpreta a resposta do FPGA como:

\begin{itemize}
  \item \codigo{word0}: sequência;
  \item \codigo{word1}: ação;
  \item \codigo{word2}: melhor compra em escala $10^4$;
  \item \codigo{word3}: quantidade na melhor compra;
  \item \codigo{word4}: melhor venda em escala $10^4$;
  \item \codigo{word5}: quantidade na melhor venda;
  \item \codigo{word6}: \textit{spread};
  \item \codigo{word7}: desbalanceamento assinado.
\end{itemize}

\section{Checklist de alteração do contrato}\label{sec:apendiceChecklistContrato}

Quando o contrato ARM--FPGA for alterado, os seguintes itens devem ser revisados em conjunto:

\begin{itemize}
  \item mapa de registradores no VHDL e na documentação;
  \item largura \codigo{G_SLOT_WORDS} no VHDL e \codigo{kFrameWords} no C++;
  \item cálculo da base RX no software, nos bancos de teste e no texto;
  \item ordem das palavras em \codigo{fast_receiver.cpp};
  \item decodificação das palavras em \codigo{order_book_core.vhd};
  \item montagem da resposta em \codigo{generated_strategy_core.vhd};
  \item testes C++ e VHDL que validam fila cheia, \textit{wrap-around} e resposta;
  \item comandos e exemplos de execução no README e neste apêndice.
\end{itemize}

Esse checklist existe porque mudanças pequenas em interfaces binárias podem causar erros difíceis de diagnosticar. Em particular, uma alteração da largura da posição muda a base RX; uma alteração da ordem das palavras pode fazer preço ser interpretado como quantidade; e uma alteração de evento sem atualização do livro pode gerar estados de topo inconsistentes.
